\section{Study Areas and Dataset}
\label{sec:data}

\subsection{Geographic and Climatic Characteristics of the Four Mountain Systems}
\label{sec:data:regions}

To rigorously assess the generalizability of deep change detection models across diverse mountain environments, this study leverages the open-access \textbf{AvalCD} benchmark \citep{avalcd_dataset, gatti2026avalanche}. The benchmark spans four prominent, geographically disparate mountain ranges across Europe, Central Asia, and the Arctic (Table~\ref{tab:dataset_summary}):

\begin{table*}[t]
\caption{Summary of the AvalCD benchmark across four global mountain systems and seven avalanche events.}
\label{tab:dataset_summary}
\centering
\footnotesize
\setlength{\tabcolsep}{6pt}
\begin{tabular*}{\textwidth}{@{\extracolsep{\fill}}l@{\hspace{4pt}}l@{\hspace{4pt}}l@{\hspace{4pt}}c@{\hspace{4pt}}r@{\hspace{4pt}}r@{\hspace{4pt}}r@{}}
\toprule
\textbf{Mountain System} & \textbf{Event Name} & \textbf{Region / Country} & \textbf{Elevation (m)} & \textbf{Grid Size (px)} & \textbf{Polygons} & \textbf{Debris (px)} \\
\midrule
\multirow{3}{*}{European Alps} & Livigno\_20240403 & Lombardy, Italy & 1,250--3,890 & $2462 \times 2091$ & 231 & 32,606 \\
 & Livigno\_20250129 & Lombardy, Italy & 1,310--3,890 & $2388 \times 1629$ & 207 & 11,438 \\
 & Livigno\_20250318 & Lombardy, Italy & 1,210--3,890 & $3069 \times 2814$ & 79 & 7,392 \\
\midrule
\multirow{2}{*}{Greenland Arctic} & Nuuk\_20160413 & Southwest Greenland & 0--1,620 & $2387 \times 4619$ & 117 & 64,822 \\
 & Nuuk\_20210411 & Southwest Greenland & 0--1,740 & $2836 \times 4931$ & 129 & 81,595 \\
\midrule
Pamir Mountains & Pish\_20230221 & Shughnon, Tajikistan & 2,100--5,420 & $1894 \times 1896$ & 113 & 22,579 \\
\midrule
Scandinavian Arctic & Tromso\_20241220 & Troms, Norway & 0--1,830 & $1691 \times 1454$ & 117 & 33,972 \\
\midrule
\textbf{Total / Benchmark} & \textbf{7 Events} & \textbf{4 Ranges} & \textbf{0--5,420} & \textbf{48.7M px} & \textbf{993} & \textbf{254,404} \\
\bottomrule
\end{tabular*}
\end{table*}

\begin{enumerate}
    \item \textbf{European Alps (Livigno, Italy):} Located in the central Italian Alps near the Swiss border ($46^\circ32'\text{N}$, $10^\circ08'\text{E}$), the Livigno region represents a temperate alpine mountain system characterized by steep glacial troughs, U-shaped valleys, and elevations ranging from 1,250~m to over 3,890~m. The region features intensive human recreation, ski infrastructure, and transportation corridors. The dataset encompasses three distinct avalanche cycles: a spring slab cycle in April 2024, a mid-winter powder cycle in January 2025, and a complex wet-snow avalanche cycle in March 2025.
    \item \textbf{Pamir Mountains (Pish, Tajikistan):} The Pamir knot in Central Asia ($37^\circ35'\text{N}$, $71^\circ39'\text{E}$) represents an extreme high-altitude, continental mountain regime characterized by deep arid valleys, narrow gorges, and rugged peaks exceeding 5,400~m in relief. In February 2023, extreme anomalous snowfall triggered a catastrophic avalanche cycle across the Shughnon District, resulting in severe infrastructure destruction and humanitarian emergencies. The snowpack is continental and facetted, with avalanche paths spanning massive vertical descents ($>2,000$~m).
    \item \textbf{Scandinavian Arctic (Troms{\o}, Norway):} Situated at $69^\circ39'\text{N}$, $18^\circ57'\text{E}$ in Northern Norway, this region features coastal maritime Arctic mountains rising directly from fjords to rugged alpine massifs. The snowpack is heavily influenced by cyclonic storms, freezing rain, and high winds, producing complex slab avalanches and wet debris. The Troms{\o} inventory provides detailed polygon annotations categorized according to the European Avalanche Warning Services (EAWS) destructive size scale (classes D1 through D4).
    \item \textbf{Greenland Sub-Polar Alpine (Nuuk, Greenland):} Surrounding the capital region of Nuuk ($64^\circ11'\text{N}$, $51^\circ44'\text{W}$), this coastal sub-polar environment consists of rugged coastal mountains, deep fjords, and permafrost terrain. Two events (April 2016 and April 2021) capture widespread avalanche releases in a polar maritime regime under intense wind-drift conditions.
\end{enumerate}

\subsection{Remote Sensing and Topographic Data Modalities}
\label{sec:data:modalities}

For each of the seven events, the benchmark provides co-registered, multi-modal geospatial rasters:
\begin{itemize}
    \item \textbf{Sentinel-1 C-band SAR Imagery:} Dual-polarization intensity data (co-polarization VV and cross-polarization VH) acquired in Interferometric Wide (IW) swath mode. For each event, a pre-event reference acquisition and a post-event acquisition are preprocessed via the ESA SNAP toolbox, applying orbit state vector updates, border noise removal, calibration to backscatter coefficient $\sigma^0$ (in decibels, dB), and Range-Doppler terrain correction at 10~m spatial resolution.
    \item \textbf{Local Incidence Angle (LIA):} Exact per-pixel local incidence angle rasters computed from satellite orbit geometry and the digital elevation model, quantifying the angle between the surface normal and the incoming radar wave.
    \item \textbf{Copernicus DEM (GLO-30):} High-quality 30~m global digital elevation model resampled to the 10~m Sentinel-1 SAR grid, providing continuous elevation measurements across all alpine catchments.
    \item \textbf{Terrain Slope and Aspect:} Derived terrain slope ($\theta_{\text{slope}} \in [0^\circ, 90^\circ]$) and aspect ($\alpha \in [0^\circ, 360^\circ]$) rasters extracted using 2D spatial gradients from the Copernicus DEM.
    \item \textbf{Ground Truth Validation Inventories:} Manual expert delineations combining optical satellite confirmation, field reports, and local avalanche bulletin records \citep{avalcd_dataset}, delivered both as raster masks ($\{0, 1\}$) and vector GeoPackages containing exact polygon boundaries and metadata.
\end{itemize}

\subsection{Partitioning for Cross-System Generalization}
\label{sec:data:split}

To limit spatial autocorrelation and characterize cross-system transfer, we partition the dataset at the regional mountain-system level. The European Alps (Livigno 2024 and Jan 2025) and Greenland (Nuuk 2016) form the training cohort (555 polygons). Model selection and threshold calibration use only Livigno March 2025 and Nuuk 2021 (208 polygons). Troms{\o} (117 polygons) and Pamir (113 polygons) are held out from training and selection. Table~\ref{tab:operational_triage} uses the finite-raster valid masks: 2,311,510 pixels in Troms{\o} and 3,503,382 in Pamir, equivalent to $231.151$ and $350.338~\text{km}^2$ on the nominal 10 m grid. At the local affine pixel areas ($85.15$ and $76.34~\text{m}^2/\text{px}$), the valid footprints are $196.820$ and $267.435~\text{km}^2$. Numerator and denominator scale together, so TopoRadar-Net's false-alarm densities are $0.38$ and $0.48~\text{ha/km}^2$ under either area convention.
